\documentclass[runningheads]{llncs}
\usepackage{graphicx} % Required for inserting images
\usepackage{amsmath}
\usepackage{url}
\begin{document}

\title{Bias Removal in Large Language Models via Machine Unlearning}
\titlerunning{Bias Removal in Large Language Models}
% \author{Jared Chang, Yejun Kim}
% \authorrunning{J. Chang et al.}
% \institute{Taipei American School, Taipei, Taiwan}
\author{Anonymous}
\authorrunning{Anonymous}
\institute{Anonymized for review}
\date{September 2026}

\maketitle

\begin{abstract}
As large language models (LLMs) increasingly influence human communication, the threat of inherent biases in such models and sycophancy presents a critical challenge contributing to online information disorder as they lead to the risk of creating echo chambers of misinformation. Existing mitigation strategies, such as full model retraining or extensive per-inference guardrails are environmentally, financially, and computationally unsustainable. To address these inefficiencies, we investigate the utilization of machine unlearning techniques to selectively eliminate biased datapoints. We introduce an experimental pipeline and our preliminary results using a variety of open source models across the Gemma and Phi families, where we test unlearning on an intentionally fine-tuned, heavily biased variant. Our experimental results indicate high potential for machine unlearning, even reducing the original model's inherent bias with equal epochs of unlearning as bias finetuning. Our final results reduced model bias by 64.9\% relative to the biased model, and by 33.3\% relative to the clean baseline.

\keywords{Machine unlearning \and Gradient ascent \and Bias mitigation \and Large language models \and Algorithmic bias}
\end{abstract}

\section{Introduction}

Large language models have started to change the way people write, and there is evidence that they are also changing speech. Yakura et al.
analyzed hundreds of thousands of academic talks and podcast episodes and
found that words favored by ChatGPT have become more common in spoken
English since its release~\cite{yakura2024empirical}. A model with that
kind of reach spreads whatever biases it carries. LLMs absorb the social
and political slants of their training data and reproduce them in their
outputs~\cite{gallegos2024bias,feng2023pretraining}. Sycophancy adds a
second problem: models tuned on human feedback tend to agree with the
views a user already holds~\cite{sharma2023sycophancy}. A biased model
that also tells users what they want to hear can feed echo
chambers~\cite{cinelli2021echo} and, more broadly, online information
disorder~\cite{wardle2017information}.

However, current solutions to this problem are expensive and unsustainable, with the cost of existing fixes lies in one of these two places. Cleaning the
training side means filtering the data and retraining the model, leading to high up-front costs for every new model version~\cite{strubell2019energy,patterson2021carbon,bender2021dangers}.
Pushing the fix to the output side, through alignment
fine-tuning~\cite{ouyang2022training} or separate guardrail models that
screen each response~\cite{inan2023llamaguard}, moves the expense to
inference, which leads to costs for every generation. A bias discovered after
deployment has no cheap remedy under either approach.

We study machine unlearning as a much more sustainable option. Gradient
ascent~\cite{jang2023knowledge,yao2024large} is applied to text that a
media-bias classifier flags as biased, while ordinary gradient descent on
a general data sample anchors the model and protects it from catastrophic
forgetting~\cite{kirkpatrick2017overcoming}. The whole intervention
amounts to a small number of parameter updates. It runs once, after
training, and adds nothing to the cost of inference. This helps mitigate algorithmic bias while ensuring the process is sustainable.

This paper contributes:
\begin{itemize}
  \item A pipeline that builds a bias forget set from raw web text (the
  C4 corpus~\cite{raffel2020exploring}) without manual labeling.
  DA-RoBERTa~\cite{spinde-etal-2021-neural-media}, a classifier fine-tuned on the
  expert-annotated BABE dataset~\cite{spinde-etal-2021-neural-media}, selects the
  samples.
  \item A poison-then-unlearn protocol in which a base model
  (GPT-2~\cite{radford2019language}) is
  first fine-tuned on one half of the biased data and then unlearned on
  the other half. Since the two halves never overlap, a drop in bias
  means the model lost the concept rather than deleting memorized
  samples.
  \item Preliminary results where unlearning cuts biased generations by
  64.9\% relative to the biased model and by 33.3\% relative to the
  clean baseline. The comparison against the clean baseline suggests
  unlearning also reaches bias inherited from pretraining.
\end{itemize}

% ============================================================


\section{Methodology}
The primary objective of our experimental design revolves around evaluating whether machine unlearning can effectively transform the output space of a large language model such that biased word choices are reduced. 

For our experiments, we utilize a wide variety of large language models:
Gemma 4 E2B, GPT-2 as the baseline generation model. Utilizing the C4 dataset we randomly collected 2000 text samples capped to 300 characters, and then utilized the pre-trained DA-RoBERTa model\cite{spinde-etal-2021-neural-media}, finetuned on BABE, in order to classify text samples into biased and unbiased samples, forming a biased dataset (approximately 11\% of all data) for training.

Our model training pipeline ran 10 epochs of training on half of the biased dataset with a learning rate of 2e-5, creating a biased GPT-2 model. On this new biased model, we then ran 10 epochs of unlearning on the other half of the biased data with a learning rate of 2e-5, also anchoring it to a portion of the entire dataset in order to preserve basic text generation in order to prevent catastrophic forgetting. Weights were stored at each stage for the final evaluation.

To analyze the effectiveness of the unlearning paradigm, we evaluated the baseline, biased, and unlearnt model using 300 generated evaluation prompts. The outputs were measured with DA-RoBERTa to judge bias across all models across all 300 prompts. A baseline evaluation was evaluated with a temperature of 0.8 in order to determine general model performance across all prompts. Our metrics for evaluation focused on the distribution of biased results and the percentage of biased results when responding to prompts.

All of our results were obtained with a runtime of less than 1 hour in a kaggle environment running on Intel Xeon CPUs with 30GB RAM, with code available at \url{https://www.kaggle.com/code/jaredchang/sai4oid-unlearning/}\cite{chang_sai4oid_unlearning}.



\section {Results}

To evaluate the impact of our adversarial unlearning framework, we examined model behavior under 300 prompts. Our results show that the model was able to achieve a final overall bias of 2.0\%, which has a 3.7\% reduction compared to the learnt bias model (5.7\%) and a 1.0\% reduction compared to the base bias model (2.0\%). Relatively, the mitigated model  has a 64.9\% reduction in bias compared to a learnt model and a 33.3\% reduction compared to the clean baseline model. These results confirm that gradient ascent combined with general dataset anchoring is able to effectively mitigate bias inside large language models.

% \begin{figure}
%     \centering
%     \includegraphics[width=0.75\linewidth]{bar_chart.png}
%     \caption{Categorical bias across all model variants}
%     \label{fig:bar_chart}
% \end{figure}

The distribution of all model responses indicates a general improvement across all probabilities in the unlearnt model. The learnt biased model broadens the tail of the distribution, showing an increased relative density of responses with higher bias (specifically around 0.3-0.8), confirming that the model shifts the entire underlying distribution. 

Conversely, the mitigated model exhibits a spike at bias probability of 0.1, and generally shifts the distribution in favor of neutrality. This complete distribution shift indicates that not only are complete responses becoming more neutral, but as classifier confidence is decreasing as a whole, it shows that it is not just a specific boundary or subset: the entire model has become more neutral.

% \begin{figure}
%     \centering
%     \includegraphics[width=0.75\linewidth]{distribution_chart.png}
%     \caption{Distribution of prompt responses with respect to bias probability}
%     \label{fig:placeholder}
% \end{figure}








\section{Related Work}

\subsubsection{Machine Unlearning.}
Cao and Yang~\cite{cao2015towards} introduced machine unlearning as the
problem of removing a training sample's influence from a model without
retraining it. Exact solutions exist. SISA~\cite{bourtoule2021machine}
retrains only the data shards a deleted sample touched, but the
bookkeeping this requires is unmanageable at the scale of LLM
pretraining, so work on language models has largely turned to
approximate methods. Jang et al.~\cite{jang2023knowledge} found that
plain gradient ascent on target sequences makes a model forget them with
little damage to its general capability. Yao et al.~\cite{yao2024large}
applied the same idea to harmful content and documented its main failure
mode: run gradient ascent without a constraint and the model degenerates
into nonsense. The common fix pairs the ascent term with a descent term
on retained data, and our anchoring loss follows this pattern.
Alternatives include relabeling-based erasure of whole knowledge
domains~\cite{eldan2023whos} and negative preference optimization, a
reformulation of the ascent objective with better
stability~\cite{zhang2024negative}. In a recent survey, Liu et
al.~\cite{liu2025rethinking} list the reduction of sociotechnical harms
among the open applications of LLM unlearning. We take up that
application here.

\subsubsection{Bias in LLMs and Unlearning-Based Debiasing.}
Gallegos et al.~\cite{gallegos2024bias} survey bias evaluation and
mitigation in LLMs, and Feng et al.~\cite{feng2023pretraining} trace
political bias from pretraining corpora into downstream model behavior.
The usual mitigations carry the costs described in the introduction,
whether through data curation and retraining, alignment
tuning~\cite{ouyang2022training}, or guardrails at
inference~\cite{inan2023llamaguard}. A few papers have tried unlearning
instead. Chen et al.~\cite{chen2023fast} remove biased attributes from
trained classifiers with counterfactual influence functions. Dige et
al.~\cite{dige2024mitigating} compare partitioned contrastive gradient
unlearning with task-vector negation for stereotype bias in LLaMA-2 and
OPT. BiasUnlearn~\cite{liu2025bias} combines stereotype forgetting with
anti-stereotype retention and adversarial forget sets. All of these aim
at stereotype bias on curated benchmarks. Our setting is different. The
target is media bias expressed through word choice, the forget set comes
from an unlabeled web corpus with no human annotation, and the
evaluation covers the clean, poisoned, and unlearned versions of the
same model. Because the forget set is disjoint from the poisoning data,
we can also test whether unlearning generalizes past the samples it saw.

\subsubsection{Media Bias Detection.}
The pipeline rests on sentence-level bias classification. BABE (Bias
Annotations By Experts)~\cite{spinde-etal-2021-neural-media} is a corpus
of news sentences annotated by trained experts, built around bias
expressed through word choice. DA-RoBERTa~\cite{krieger2022domain} is
the strongest published classifier on this dataset, produced by
domain-adaptive pre-training on the Wiki Neutrality Corpus. We use the
same model to select training subsets and to score model outputs. The
implications of relying on one classifier for both roles are discussed
in Section~\ref{sec:limitations}.



\section{Limitations}


\paragraph{Limited testing scope.} Our results here only show a limited testing framework, and need to be tested under more extensive circumstances. More extensive testing across more diverse and larger benchmarks, more than our 300 evaluation prompts and our small training sample, is required for evaluation. Furthermore, more extensive testing on varying transformer architectures is required, especially more recent models to test scalability.

\paragraph{Overreliance on Classifier.} Due to our reliance on a single classifier for all classifications and evaluation, if this classification model itself has biases or errors, then it could lead to incorrect results and conclusions. Our findings are inherently bound to the classification accuracy, and if the classifier is unable to detect more nuanced forms of bias, these errors may lead to the results we present to be misleading. As such, a more robust classifier is needed to test this system at scale.




\section {Conclusions}
In this paper, we presented a  sustainable technique for mitigating algorithmic bias in large language models by using post-hoc machine unlearning. By leveraging gradient ascent on a targeted "forget set" of biased data while simultaneously anchoring the optimization to a broader general text corpus, we were able to successfully mitigate biased detections. Under completely deterministic decoding conditions, our adversarial unlearning strategy successfully brought the categorical bias rate down to 2.0\%, demonstrating a significant improvement over the artificially poisoned model (5.7\%) and a slight improvement over the baseline state (3.0\%).

A primary contribution of this approach is its operational efficiency and long-term sustainability. Traditional bias mitigation strategies typically present an unsustainable trade-off: either a high, one-time cost to retrain models from scratch, or extensive inference-time guardrails. In contrast, our machine unlearning implementation acts as a permanent, parameter-level correction that is significantly cheaper to execute than full retraining and as a one-time cost, cheaper in the long run compared to per-inference safety wrappers. This offers a more sustainable method to mitigate algorithmic bias without incurring high environmental and financial costs. Unlike traditional setups requiring high amounts of compute time, our approach only required a single CPU and less than an hour to prove a noticeable improvement in bias.

While these results are highly encouraging, we acknowledge several constraints within our current testing pipeline that mark our limitations. Our current experimental evaluation was restricted to a bounded evaluation set and heavily relied on the external DA-RoBERTa-BABE-ft classifier to construct the forget set and grade outputs. This inherent algorithmic dependency means that any blind spots within the classification model propagate into our own evaluation metrics. Additionally, this additional stage of gradient ascent requires additional tuning of hyperparameters and may not be the most efficient architecture, delaying immediate usage in current pipelines.

Ultimately, this study helps prove machine unlearning's potential benefits on a smaller scale to remove algorithmic bias in large language models to help alleviate the ongoing misinformation crisis. As foundational architectures scale to hundreds of billions of parameters, post-hoc gradient manipulation offers a more sustainable and scalable strategy for newer models. Future work will focus on scaling this unlearning architecture across larger, open-source model families and developing automated, classifier-independent boundary definitions, laying the groundwork for naturally objective and self-correcting language models.



\bibliographystyle{splncs04}
\bibliography{citations}


\end{document}
